\section{记录判断的配对统计}
\label{sec:statistics}

\subsection{结果定义：记录通过}
\label{subsec:pass-coding}
沿着选定的候选修改，较晚端点记录通过的频率比早期端点增加了多少？较大的历史档案允许我们把相应的前后记录配对，计算这一变化。

本节使用按第~\ref{subsec:revision-rule}~节活动规则选出的全部 465 条候选修改边，涉及 250 个任务；配对分析不在首次通过处停止。每条边比较同一作业题的两个版本，一个任务可以贡献多次比较。因此，前后记录是配对观察，而任务是归组重复观察的单位。

先将记录标签转换成数值结果。记录写明\code{pass}，即通过时，记为 $A=1$。对于明确记录为\code{fail}、\code{inconclusive}、\code{partial}，即失败、无法定论、部分完成的情形，记为 $A=0$。因此，$A$ 只回答“是否记录了通过”。零不表示数学内容一定错误；尤其是无法定论的审查，并未认定数学失败。分析脚本会拒绝缺失或无法识别的标签，而不是悄悄将其补成零。表~\ref{tab:labels} 还保留了原来的四个类别。

对一条边，用较晚值减去较早值。$0\to1$ 贡献 $+1$，$1\to0$ 贡献 $-1$，两端二元值相同则贡献 $0$。这些量都没有直接衡量语义正确性；所计算的结果始终是记录中的审查决定。

\subsection{任务等权与边等权平均变化}
\label{subsec:task-weighting}
考虑表~\ref{tab:task-weights}~中的说明性例子。任务甲在一次选定修改后通过，任务乙直到第三次才通过。两个任务最终都通过了，但记录的边数不同。

\begin{table}[htbp]
\centering\small
\caption{仅用于说明的例子：不同任务贡献的边数可以不同。表中数值表示是否记录通过，不是数学质量分数。}
\label{tab:task-weights}
\begin{tabularx}{\linewidth}{l c Y c}
\toprule
任务 & 边数 & 较早值 $\to$ 较晚值 & 边差值的平均数 \\
\midrule
甲 & 1 & $0\to1$ & $1$ \\
乙 & 3 & $0\to0$、$0\to0$、$0\to1$ & $1/3$ \\
\bottomrule
\end{tabularx}
\end{table}

若每个任务只占一份权重，平均数是 $(1+1/3)/2=2/3$。若每条边各占一份权重，平均数是 $(1+0+0+1)/4=1/2$。两者不同，是因为后一种算法给予任务乙的权重是任务甲的三倍。两种算术运算都没有错；它们分别回答“各任务自身的平均变化，再按任务平均是多少”和“所有已记录边的平均变化是多少”。它们都不是最终通过的任务比例；在这个例子中，最终通过比例是 $2/2$。

现在为实际数据中的这些对象逐一命名。用 $S$ 表示纳入的任务数，主分析中 $S=250$。下标 $k$ 表示任务编号，从 $1$ 到 $S$；$n_k$ 是任务 $k$ 的选定边数。在任务 $k$ 内，下标 $e$ 从 $1$ 到 $n_k$，用于标记其中的边。$A_{ke,0}$ 和 $A_{ke,1}$ 分别表示这条边的较早和较晚二元值。定义
\begin{equation}
\begin{split}
d_k&=\frac{1}{n_k}\sum_{e=1}^{n_k}(A_{ke,1}-A_{ke,0}),\\
\widehat\mu_{\mathrm{task}}&=\frac{1}{S}\sum_{k=1}^{S}d_k,
\qquad
\widehat\mu_{\mathrm{edge}}=\frac{\sum_{k=1}^{S}n_kd_k}{\sum_{k=1}^{S}n_k}.
\end{split}
\label{eq:paired-estimands}
\end{equation}
其中，$d_k$ 是一个任务内部的平均边差值，$\widehat\mu_{\mathrm{task}}$ 是任务等权平均数，$\widehat\mu_{\mathrm{edge}}$ 是边等权平均数。这三个量都是零到一尺度上的无量纲差值，乘以 100 才是百分点。$d_k$ 也等于该任务的较晚端点均值减去较早端点均值；后面的统计计算，配对的正是每个任务的这两个均值。

在完整保留集合中，较早端点出现 118 次通过，较晚端点出现 328 次通过。因此
\[
\widehat\mu_{\mathrm{edge}}=\frac{328-118}{465}
=\frac{210}{465}\approx0.451613,
\]
即增加 45.16 个百分点。若先算每个任务的差值，再对 250 个任务平均，则得到 $\widehat\mu_{\mathrm{task}}=0.570909$，即增加 57.09 个百分点。在引入任何概率模型之前，这两个数就已经是对指定档案与纳入规则的精确描述；展示的小数经过了舍入。

\subsection{为什么 465 条边不是 465 个独立观察}
\label{subsec:edge-dependence}
考虑记录值为 $0\to0\to1$ 的一条链。两个相邻差值共用了中间记录：它既是前一条边的较晚端点，也是后一条边的较早端点。更一般地，若一条链有 $m$ 条边，各端点数值依次为 $a_0,a_1,\ldots,a_m$，则
\[
(a_1-a_0)+(a_2-a_1)+\cdots+(a_m-a_{m-1})=a_m-a_0.
\]
所有中间值都互相抵消。因此，一条长链可以提供很多行数据，却不一定提供同样多的独立净变化观测。

保留的图中有 930 次端点出现，但只有 783 条不同的端点记录；其中 147 条记录出现在两条边中。其 318 条极大链使用全部 465 条边构成，与 367 条高置信边所组成的 262 个过程片段不同；这些链满足上述抵消恒等式，总净差值为 210。共同的任务难度、反馈和审查条件还可能带来额外依赖。103 个任务具有多条边，单个任务最多有 26 条边；边数最多的五个任务占全部边的 14.41\%。这里的图链仅用于诊断这批边的结构，并未替换后文失败起点过程的定义。

所以，我们先将每个任务概括为配对差值 $d_k$，再计算 t 统计量；计算参考区间时也以整个任务为单位重抽样。这样保留了任务\emph{内部}的依赖结构，但不能证明不同任务彼此独立：它们仍可能共用审查者、批次、数学依赖或不断变化的规则。

\subsection{配对 t 计算及其参照假设}
\label{subsec:paired-t-explanation}
\subsubsection{平均差与标准误}
本次有条件计算使用 250 个任务层面的配对差值，参照模型要求任务独立，并遵循同一种记录选择机制。

令 $D$ 为参照总体中的任务差值，均值为 $\mu_D=E(D)$，标准差为 $\sigma_d$。对于观察到的 $d_1,\ldots,d_S$，
\[
s_d=\sqrt{\frac{1}{S-1}\sum_{k=1}^{S}(d_k-\widehat\mu_{\mathrm{task}})^2}.
\]
观察样本的标准差为 $s_d=0.430341$，即 43.03 个百分点。

在任务差值相互独立、共同方差为 $\sigma_d^2$ 的条件下，
\[
\operatorname{Var}\!\left(\frac{1}{S}\sum_{k=1}^{S}D_k\right)
=\frac{1}{S^2}\sum_{k=1}^{S}\sigma_d^2
=\frac{\sigma_d^2}{S}.
\]
估计标准误为
\[
\widehat{\mathrm{SE}}=\frac{s_d}{\sqrt S}
\approx\frac{0.430341}{\sqrt{250}}\approx0.027217.
\]
在该任务模型下，它是 2.72 个百分点。跨任务依赖会引入协方差项，因此直接除以 $\sqrt{250}$ 未必能描述实际档案中均值的不确定性。

\subsubsection{零假设与观察统计量}
配对零假设为指定任务总体中的 $H_0:\mu_D=0$。保存的双侧计算为\cite{nistpaired,scipyttest}
\begin{equation}
t_{\mathrm{obs}}=\frac{\widehat\mu_{\mathrm{task}}-0}{s_d/\sqrt S}
\approx\frac{0.5709091575}{0.0272171622}\approx20.9761.
\label{eq:paired-t-worked}
\end{equation}
配对的两项是每个任务的较晚端点均值和同一任务的较早端点均值。把前后集合当成独立样本会丢掉这层配对；本节的模型针对的是两者之差。

\subsubsection{参照分布与适用范围}
对于独立同分布、具有正方差的正态任务差值，零假设下的参照分布为自由度 $S-1=249$ 的 t 分布。

实际任务差值有上下界，而且只取离散值，因此并不严格满足正态差值模型。对于这 250 个任务，报告仅在额外的独立任务模型下，将 Student 计算用作探索性的大样本参照。尤其需要注意，任务数量较大，既不能认证任务独立性，也不能保证极小尾概率的数值校准准确。

令 $Z\sim t_{249}$，名义双侧 p 值为
\[
p_{\mathrm{ref}}=\Pr\{|Z|\geq |20.9761|\}
\approx5.87\times10^{-57}.
\]
这是参照模型下的名义尾概率；该模型是否适用于这批选定档案尚未得到确认。数值很小并不能建立数学改善。

\subsection{任务整群重抽样与参考区间}
\label{subsec:cluster-resampling}
为了考察所含任务构成的变化，保存的分析采用以\emph{整个任务}为抽样群的自助法\cite{scipybootstrap}，每次抽取都携带该任务全部选定配对边。

每次重复，都从已有的 $S$ 个任务编号中，等概率、独立地抽取 $S$ 次。被抽中的任务会带上全部边和每条边的两个端点。任务等权均值由抽中任务的差值重新平均。边等权均值则用抽中任务的总净变化除以总边数，所以其分母在不同重复之间可能变化。

保存的分析执行了 20,000 次重复。将所得数值排序，取第 2.5 和第 97.5 百分位数，得到报告中的中央 95\% 百分位区间。完整数据的任务等权区间为 $[0.518357,0.624482]$，即 $[51.84,62.45]$ 个百分点；边等权区间为 $[38.63,52.12]$ 个百分点。复算使用 NumPy PCG64、种子序列 $[20260906,r]$ 和线性百分位插值，其中 $r=0,1,2$ 标记三条分析路线。

这里称为\emph{参考}区间，是因为通常的 95\% 覆盖率解释需要任务彼此独立、足够可比，并由同一种历史及纳入机制产生。这批回顾性档案尚未确认这些假设。给定 250 个观察差值，均值 $0.570909$ 就是固定的；区间描述的是假想任务构成变化。按整任务重抽样保留了任务内部依赖，但不能修复选择偏差、跨任务依赖或缺少语义结果测量的问题。

这与\emph{逻辑补全界}不同。例如，设较早分数为零，较晚分数尚未测量，但只能是零或一。差值便可能为零或一，紧确界为 $[0,1]$。这些可能性没有指定概率，也不带有 95\% 覆盖率的含义；测得缺失分数后，歧义就可能消失。任务自助法则使用已知分数，改变假想样本的构成。

\subsection{三种选取规则下的结果}
\label{subsec:paired-results}
表~\ref{tab:paired-effects} 和表~\ref{tab:paired-tests} 列出主分析及两条敏感性路线。敏感性路线改变某项纳入规则，展示这个变化后的问题会得到什么答案。三条路线是在上一版修订期间、数值估计之前声明的，但当时已经了解这批档案，因此这是回顾性探索，不是前瞻预注册。三项未经多重比较校正的双侧参考检验均予报告；没有按 p 值大小择取路线，逐项区间也不提供三条路线同时成立的置信保证。

\begin{table}[htbp]
\centering\small
\caption{记录通过的变化，单位为百分点。方括号内为以任务模型为前提的 95\% 任务整群百分位参考区间。改变路线，也改变了纳入的边和任务。}
\label{tab:paired-effects}
\begin{tabularx}{\linewidth}{Y r r r}
\toprule
路线 & 边数 / 任务数 & 任务等权变化 & 边等权变化 \\
\midrule
全部候选修改边 & 465 / 250 & 57.09 [51.84, 62.45] & 45.16 [38.63, 52.12] \\
两端均为通过或失败 & 440 / 236 & 55.73 [50.15, 61.21] & 44.09 [37.36, 51.18] \\
高置信候选修改边 & 367 / 207 & 76.93 [71.89, 81.73] & 59.95 [50.44, 70.10] \\
\bottomrule
\end{tabularx}
\end{table}

\begin{table}[htbp]
\centering\small
\caption{对任务层面均值进行的配对计算。前后百分比均为任务等权。名义 p 值使用尚未验证的 Student 参照模型，不检验语义改善。}
\label{tab:paired-tests}
\begin{tabularx}{\linewidth}{Y c c c}
\toprule
路线 & 较早 $\to$ 较晚 & $t$（df） & 双侧 $p_{\mathrm{ref}}$ \\
\midrule
全部候选修改边 & 29.16\% $\to$ 86.25\% & 20.9761 (249) & $5.87\times10^{-57}$ \\
两端均为通过或失败 & 31.59\% $\to$ 87.32\% & 19.7176 (235) & $1.00\times10^{-51}$ \\
高置信候选修改边 & 6.92\% $\to$ 83.85\% & 30.5239 (206) & $2.21\times10^{-78}$ \\
\bottomrule
\end{tabularx}
\end{table}

完整集合中，182 个任务的差值为正，67 个为零，1 个为负。数据中确实存在变化，因此标准差非零，t 统计量可以计算。任务等权与边等权均值之间的差距也有实际意义：选定历史较长的任务在第二种均值中权重更大，而它们合起来的每边净变化较小。

\begin{table}[htbp]
\centering\small
\caption{全部 465 条边的原始审查标签。行表示较早端点，列表示较晚端点。即使“无法定论”和“部分完成”都没有编码为通过，它们仍在表中单独呈现。}
\label{tab:labels}
\begin{tabular}{lrrrr}
\toprule
较早 / 较晚 & 通过 & 失败 & 无法定论 & 部分完成 \\
\midrule
通过 & 105 & 12 & 1 & 0 \\
失败 & 206 & 117 & 5 & 1 \\
无法定论 & 16 & 0 & 1 & 0 \\
部分完成 & 1 & 0 & 0 & 0 \\
\bottomrule
\end{tabular}
\end{table}

四分类表解释了编码为什么重要。25 条边至少有一端为无法定论或部分完成；其中 17 条从这种结果转为通过，1 条从通过转为这种结果。因此主均值还包含了“是否记录了明确通过”的变化。限制两端都只能为通过或失败，会去掉这些边，但也改变纳入的任务。这是另一个经过选择的目标，而不是给原样本换了一个更好的名字。

高置信路线沿用的是候选修改边的置信分类；“高置信”不是数学修复已获独立认证。该分类规则使用此前的结构化问题记录。在主估计之后增加的一项描述检查显示，两组的起点构成差异很大：高置信子集在 367 条边中有 21 个较早通过端点；与之互补的 98 条中置信边，开始时有 97 个通过，结束时有 87 个。这个互补集合涉及 81 个任务，任务等权变化为 $-11.73$ 个百分点，边等权变化为 $-10.20$ 个百分点。本项估计后检查没有附加新的检验或区间。起点几乎全部通过的组，向上增长的空间很小；起点大多未通过的组，空间则大得多。因此，高置信组变化更大，并不能证明其修改更有效。

其他解释也没有排除。45 条边的来源审查契约字段不同，2 条边的任务内容字段不同；字段一致并不证明实际审查条件完全相同。在 783 条不同端点记录中，781 条的检查范围未知，1 条明确报告未审查，1 条报告存在可审查性阻碍。本次投影没有检查独立执行。451 条边处于单独案例评估之外，因此较大统计集合中的这些边缺少对应语义评估。候选本身改善、评估条件变化和选择性继续，都可能带来记录标签变化。向均值回归也可能带来部分变化。没有合适的比较设计和独立语义结果，仅凭很小的 p 值无法区分这些解释。

\subsection{为什么零诊断差值给出未定义的 t 统计量}
\label{subsec:zero-differences}
较小的评价器比较测量另一种结果：\emph{预先规定的诊断问题中，能够给出确定答案的数量}。四格方法的基线在 $R_1$ 下检查两个候选，同时使用已经核实的标准关系；由这些关系得到的推论，也属于基线可用的信息。

对任务 $k$，令 $C_{k,4}$ 为四格方法给出的确定答案数，$C_{k,F}$ 为固定要求基线的相应答案数。每个计数介于零和六之间。定义 $D_k=C_{k,4}-C_{k,F}$，其单位是\emph{诊断问题数}，不是百分点。在两种历史评估中，各自的 14 个 $D_k$ 都是零，而且逐个答案也一致。套入相同计算步骤，得到
\[
\overline D=0,\qquad
s_D=\sqrt{\frac{1}{13}\sum_{k=1}^{14}(D_k-\overline D)^2}=0,
\qquad t_{\mathrm{obs}}=\frac{0}{0/\sqrt{14}}=\frac{0}{0}.
\]
该统计量未定义，并非得到 $t=0,p=1$ 的 t 检验。前面的 250 任务计算中存在非零变化，分母可以定义；这里没有。

精确的描述性发现仍然有用：对于这些选定历史案例，两种方法给出的答案相同。但任务自助法只能重新抽出一串零，符号方向比较也没有正差或负差可比较。这些计算不能揭示未观察任务是否会出现差异。将相同 14 个任务重复用于两种评估，或分别数它们的六个问题和四个格子，都不能制造独立任务信息。若要声称总体等效，需要明确总体、合理的不确定性模型，并事先说明等效问题中多大的差异在实际意义上可以接受。

七个构造例子相对固定基线的增益为 $[5,0,4,0,0,0,0]$。它们证明额外信息\emph{可能}有帮助；刻意安排的样例构成，不能估计这种帮助在普通任务中有多常见。夏普利分配相对直接读取四格，在全部 35 张表中都没有增加诊断可答性，这与它使用相同观测信息一致。这些是不同的比较，不应混成“所有地方的增益都为零”这一说法。

\subsection{八对编译结果：精确 McNemar 检验}
\label{subsec:mcnemar}
八个任务隔离留出候选对还支持另一项可计算的比较：较早与较晚候选在当前测试环境中是否能够编译？表~\ref{tab:compile-pairs} 给出完整配对观察。这一结果只涉及编译，不表示预期数学命题已经得到正确表达。

\begin{table}[htbp]
\centering\small
\caption{八个留出任务对的实际编译结果。提供变化方向信息的是两个非对角单元格，而不是将 16 个端点分开后视为独立观测。}
\label{tab:compile-pairs}
\begin{tabular}{lrrr}
\toprule
较早 / 较晚 & 编译失败 & 编译通过 & 合计 \\
\midrule
编译失败 & 1 & 1 & 2 \\
编译通过 & 1 & 5 & 6 \\
\midrule
合计 & 2 & 6 & 8 \\
\bottomrule
\end{tabular}
\end{table}

五对两端均能编译，一对两端均不能；一对从失败转为成功，一对从成功转为失败。因此，前后均为八个候选中六个编译通过，配对净差为 $(1-1)/8=0$。改善的是任务\code{def\_13\_2}，退步的是任务\code{prob\_10\_1}。

精确 McNemar 计算所问的是：不一致配对的方向是否平衡\cite{nistmcnemar}？令 $b=1$ 记录失败转成功的对数，$c=1$ 记录反方向的对数，则不一致配对总数为 $m=b+c=2$。在不同不一致方向相互独立、两种方向等概率的条件零假设下，假想重复中改善方向的个数 $B$ 服从试验次数为 $m=2$、成功概率为 $1/2$ 的二项分布。采用较小尾概率的两倍、最多截为一的双侧规则，得到
\[
p=\min\{1,\,2\Pr(B\leq1)\}
=\min\{1,\,2(1/4+1/2)\}=1.
\]
这个 p 值确实有定义，但只有两个提供信息的方向。它没有给出反对上述零假设的方向性证据；也没有证明编译概率等效、修改没有效果或语义正确。留出选择经过分层且具有回顾性，任务隔离本身不能证明统计独立，也不能抹去先前暴露。这项检验只是适用于配对二元数据的一次有限计算示例，不能代替较大集合的任务层面分析。

配对差概括了选定修改的变化。若要研究这些变化的先后顺序，就需要沿路径逐步观察，直到首次记录通过或路径退出。